\documentclass{situationist-preprint}

\title{Invertebrate Sentience and Moral Status}
\msid{ETH-07-b}
\msfamily{ETH-07}
\msdiscipline{Normative and Applied Ethics}
\mstype{Companion manuscript}
\mskeywords{invertebrates; sentience; moral status; cephalopods; decapods; insects; precaution; uncertainty}

\begin{document}
\maketitle

\begin{abstract}
ETH-07-a argues that sentience is necessary and sufficient for moral status and that degree of status varies with capacity for welfare. For invertebrates, the first question is therefore empirical. We review the evidence for sentience in cephalopods, decapod crustaceans and insects, and argue that it is strong for the first, substantial for the second and inconclusive for the third. We then show how the view of ETH-07-a guides action under this uncertainty: since uncertainty about status is uncertainty about sentience, the weight owed to an invertebrate's possible suffering is its probability of sentience times what would be at stake for it if sentient, and precautions are owed in proportion.
\end{abstract}

\section{The evidence}

By Theorem~3.2 of ETH-07-a, a being has moral status if and only if it is sentient, that is, capable of valenced experience (Definition~2.2 of ETH-07-a). Whether invertebrates have moral status is therefore the question whether they are sentient, and the evidence bears on it through behavioral and neural markers: integration of noxious stimuli with other information, motivational trade-offs, wound-directed behavior, learned avoidance, responses to analgesics, and affect-like states \citep{Birch2021}.

\paragraph{Cephalopods.} Octopuses show wound-directed behavior, learned avoidance of places associated with noxious stimulation, and preference for places where they received a local anaesthetic, which together suggest an affective component of pain and not mere nociception \citep{Crook2021}. Their nervous systems are large and highly integrated, though distributed. MND-23 concludes that cephalopods are probably conscious, and the evidence for valence specifically is among the strongest outside the vertebrates.

\paragraph{Decapod crustaceans.} Hermit crabs subjected to electric shocks within their shells trade off escape from the shock against the quality of the shell they would have to leave, accepting more shock to keep a better shell \citep{Elwood2009}. Such trade-offs indicate that the noxious stimulus is weighed against other motivations in a central evaluation, which is a mark of felt aversiveness. A systematic review applying eight criteria found strong evidence of sentience in decapods as well as cephalopods \citep{Birch2021}. We regard the evidence as substantial, though less strong than for cephalopods.

\paragraph{Insects.} Bumblebees show motivational trade-offs in response to noxious heat \citep{Gibbons2022} and behavior suggestive of positive affect-like states after unexpected rewards \citep{Solvi2016}, and it has been argued that the insect midbrain performs functions analogous to those associated with consciousness in vertebrates \citep{Barron2016}. But the markers are fewer and the interpretations more contested. With MND-23, we hold that the evidence for insects remains inconclusive.

\section{Acting under uncertainty}

ETH-07-a holds that uncertainty about whether a being has moral status is uncertainty about an empirical fact, sentience, and not about the grounds of status (Section~6 of ETH-07-a). This allows the uncertainty to be handled by ordinary expected-value reasoning applied to the moral weight of welfare changes.

\begin{proposition}\label{p:expected}
Under equal consideration of interests (Definition~5.1 of ETH-07-a), if $p$ is the probability that a being is sentient and $\Delta$ the welfare change an act would cause it if it is sentient, the expected moral weight of the act's effect on the being is $p\,\Delta$.
\end{proposition}

\begin{proof}
If the being is sentient, it is a subject of welfare (Lemma~3.1 of ETH-07-a), the change $\Delta$ is a change in its welfare, and by equal consideration it carries weight $\Delta$. If it is not sentient, it is not a subject of welfare, nothing that happens to it is a change in its welfare, and the effect carries weight $0$. The expectation is $p\,\Delta + (1-p)\cdot 0$.
\end{proof}

Proposition~\ref{p:expected} has three consequences. First, cephalopods and decapods are owed substantial consideration now: with a high probability of sentience, the expected weight of severe suffering, such as that caused by boiling alive, is close to its full weight. The inclusion of both groups in the protections of the United Kingdom's Animal Welfare (Sentience) Act 2022, on the basis of the review by \citet{Birch2021}, is what the view recommends. Second, insects are owed consideration discounted by the probability of their sentience, and by the small capacity for welfare that, if sentient, they probably have (Proposition~3.3 of ETH-07-a); but since the numbers affected by farming, pesticides and land use are very large, the aggregate weight need not be small. Third, the discount is not a reduction of status: an insect that is sentient has status of the degree its capacity fixes, and the probability measures our evidence, not its standing.

This agrees with the precautionary framework of \citet{Birch2024}, on which beings for which there is a realistic, evidence-based possibility of sentience warrant proportionate precautions. Proposition~\ref{p:expected} supplies one way of fixing the proportion.

\begin{objection}
Expected-value reasoning about tiny probabilities and vast numbers leads to fanaticism: any being with a non-zero chance of sentience, in sufficient numbers, would dominate moral concern.
\end{objection}

\begin{reply}
The probabilities in question are not tiny. Insects meet several of the markers, and a probability of sentience that is inconclusive is far from negligible. Where there is no evidence at all, as for plants, ETH-07-a argues that there are no grounds for attributing valenced experience, and the probability is not one on which action should be based. The objection applies, if at all, to beings outside the class of sentience candidates, and the view does not ask us to act for them.
\end{reply}

\section{Conclusion}

Cephalopods are probably sentient, decapods substantially likely to be, and the evidence for insects is inconclusive. Since sentience suffices for moral status, the first two are owed consideration now, and insects are owed consideration discounted by the probability of their sentience but multiplied by their numbers.

\begin{thebibliography}{9}\small
\bibitem[Barron and Klein(2016)]{Barron2016} Barron, A.~B. and Klein, C. (2016). What insects can tell us about the origins of consciousness. \emph{Proceedings of the National Academy of Sciences}, 113(18), 4900--4908.
\bibitem[Birch(2024)]{Birch2024} Birch, J. (2024). \emph{The Edge of Sentience: Risk and Precaution in Humans, Other Animals, and AI}. Oxford University Press.
\bibitem[Birch et~al.(2021)]{Birch2021} Birch, J., Burn, C., Schnell, A., Browning, H. and Crump, A. (2021). \emph{Review of the Evidence of Sentience in Cephalopod Molluscs and Decapod Crustaceans}. London School of Economics and Political Science.
\bibitem[Crook(2021)]{Crook2021} Crook, R.~J. (2021). Behavioral and neurophysiological evidence suggests affective pain experience in octopus. \emph{iScience}, 24(3), 102229.
\bibitem[Elwood and Appel(2009)]{Elwood2009} Elwood, R.~W. and Appel, M. (2009). Pain experience in hermit crabs? \emph{Animal Behaviour}, 77(5), 1243--1246.
\bibitem[Gibbons et~al.(2022)]{Gibbons2022} Gibbons, M., Versace, E., Crump, A., Baran, B. and Chittka, L. (2022). Motivational trade-offs and modulation of nociception in bumblebees. \emph{Proceedings of the National Academy of Sciences}, 119(31).
\bibitem[Solvi et~al.(2016)]{Solvi2016} Solvi, C., Baciadonna, L. and Chittka, L. (2016). Unexpected rewards induce dopamine-dependent positive emotion-like state changes in bumblebees. \emph{Science}, 353(6307), 1529--1531.
\end{thebibliography}

\end{document}
